% GENERATED FILE. Edit canonical lessons, then run npm run generate:books.
% canonical-source-hash: fnv1a64:1d917eca58fb62a7
% canonical-lessons: LA-C344-maiores, LA-C344-annales, LA-C344-conditor, LA-C344-origo, LA-C344-vetustas

\chapter{Maiōrēs nostrī}
\label{ch:la-maiores-nostri}
\hlchaptermodality{344}

\begin{chapteropening}
\textbf{By the end of this chapter:} \emph{I can talk about the past and its records: the ancestors, the annals, the founder, the origin and great age.}

Complete the last lesson of chapter 344: I can talk about the past and its records: the ancestors, the annals, the founder, the origin and great age.
\end{chapteropening}

\section[maiōrēs, maiōrum]{maiōrēs, maiōrum — ancestors, forefathers}

\label{lesson:LA-C344-maiores}



\begin{quote}
Before the new one: say the Latin for a treaty, then the Latin for a triumph.
\end{quote}

\subsection*{You'll want to know: maiōrēs, maiōrum}
\textbf{maiōrēs, maiōrum} — \textquotedblleft{}ancestors, forefathers\textquotedblright{}.

\begin{quote}
Maiōrēs nostrī hoc templum aedificāvērunt.
\end{quote}

\begin{quote}
\emph{Our ancestors built this temple.}
\end{quote}

Masculine, plural: the comparative of \emph{magnus} as a noun, the 'greater' generations before us.

\begin{cousinweb}
\textbf{maiōrēs} is the plural of \emph{maior}, \textquotedblleft{}greater, elder\textquotedblright{}, the comparative of \emph{magnus}, \textquotedblleft{}great\textquotedblright{} (Lewis \& Short): the elders, the ancestors.
\end{cousinweb}

\begin{grammarlens}[title={What you've built}]
One more word for telling, explaining and describing at B1.
\end{grammarlens}

\subsection*{Your turn}
\begin{itemize}
\raggedright
  \item \emph{Say it:} \emph{maiōrēs, maiōrum}
  \item \emph{Say it:} \emph{maiōrēs, maiōrum}, once more
  \item \emph{Say it:} a sentence of your own with \emph{maiōrēs, maiōrum}
  \item \emph{Recall:} say the Latin for a treaty, then the Latin for a triumph, then say \emph{maiōrēs, maiōrum} again
\end{itemize}

\begin{tcolorbox}[breakable,colback=teal!4,colframe=teal!35!black,title={Before you move on}]
What does maiōrēs, maiōrum mean? (\textquotedblleft{}Ancestors, forefathers\textquotedblright{}.) Say it once more.
\end{tcolorbox}
\section[annālēs, annālium]{annālēs, annālium — annals, a year-by-year history}

\label{lesson:LA-C344-annales}



\begin{quote}
Before the new one: say the Latin for a triumph, then the Latin for ancestors.
\end{quote}

\subsection*{You'll want to know: annālēs, annālium}
\textbf{annālēs, annālium} — \textquotedblleft{}annals, a year-by-year history\textquotedblright{}.

\begin{quote}
In annālibus omnia bella scrīpta sunt.
\end{quote}

\begin{quote}
\emph{All the wars are written in the annals.}
\end{quote}

Masculine, plural, with \emph{librī} understood. The priests kept a record of each year; later historians wrote \emph{annālēs}.

\begin{cousinweb}
\textbf{annālēs} is built on \emph{annus}, \textquotedblleft{}year\textquotedblright{} (EtymDB): the yearly books. English \emph{annals} comes from it.
\end{cousinweb}

\begin{grammarlens}[title={What you've built}]
One more word for telling, explaining and describing at B1.
\end{grammarlens}

\subsection*{Your turn}
\begin{itemize}
\raggedright
  \item \emph{Say it:} \emph{annālēs, annālium}
  \item \emph{Say it:} \emph{annālēs, annālium}, once more
  \item \emph{Say it:} a sentence of your own with \emph{annālēs, annālium}
  \item \emph{Recall:} say the Latin for a triumph, then the Latin for ancestors, then say \emph{annālēs, annālium} again
\end{itemize}

\begin{tcolorbox}[breakable,colback=teal!4,colframe=teal!35!black,title={Before you move on}]
What does annālēs, annālium mean? (\textquotedblleft{}Annals, a year-by-year history\textquotedblright{}.) Say it once more.
\end{tcolorbox}
\section[conditor, conditōris]{conditor, conditōris — a founder}

\label{lesson:LA-C344-conditor}



\begin{quote}
Before the new one: say the Latin for ancestors, then the Latin for annals.
\end{quote}

\subsection*{You'll want to know: conditor, conditōris}
\textbf{conditor, conditōris} — \textquotedblleft{}a founder\textquotedblright{}.

\begin{quote}
Conditor urbis in forō sepultus esse dīcēbātur.
\end{quote}

\begin{quote}
\emph{The founder of the city was said to be buried in the forum.}
\end{quote}

Masculine, from \emph{condō}, 'to found' (taught), short \emph{i}. Do not confuse \emph{condītor}, 'a seasoner'.

\begin{cousinweb}
\textbf{conditor} is built on \emph{condō}, \textquotedblleft{}to put together, to found\textquotedblright{} (Lewis \& Short), with the ending \emph{-tor}.
\end{cousinweb}

\begin{grammarlens}[title={What you've built}]
One more word for telling, explaining and describing at B1.
\end{grammarlens}

\subsection*{Your turn}
\begin{itemize}
\raggedright
  \item \emph{Say it:} \emph{conditor, conditōris}
  \item \emph{Say it:} \emph{conditor, conditōris}, once more
  \item \emph{Say it:} a sentence of your own with \emph{conditor, conditōris}
  \item \emph{Recall:} say the Latin for ancestors, then the Latin for annals, then say \emph{conditor, conditōris} again
\end{itemize}

\begin{tcolorbox}[breakable,colback=teal!4,colframe=teal!35!black,title={Before you move on}]
What does conditor, conditōris mean? (\textquotedblleft{}A founder\textquotedblright{}.) Say it once more.
\end{tcolorbox}
\section[orīgō, orīginis]{orīgō, orīginis — an origin, a beginning; ancestry}

\label{lesson:LA-C344-origo}



\begin{quote}
Before the new one: say the Latin for annals, then the Latin for a founder.
\end{quote}

\subsection*{You'll want to know: orīgō, orīginis}
\textbf{orīgō, orīginis} — \textquotedblleft{}an origin, a beginning; ancestry\textquotedblright{}.

\begin{quote}
Dē orīgine urbis multae fābulae nārrantur.
\end{quote}

\begin{quote}
\emph{Many stories are told about the origin of the city.}
\end{quote}

Feminine, 3rd declension, from \emph{orior}, 'to rise'. English \emph{origin}.

\begin{cousinweb}
\textbf{orīgō} is built on \emph{orior}, \textquotedblleft{}to rise, to arise\textquotedblright{} (Lewis \& Short): the rising, the source. English \emph{origin} comes from it.
\end{cousinweb}

\begin{grammarlens}[title={What you've built}]
One more word for telling, explaining and describing at B1.
\end{grammarlens}

\subsection*{Your turn}
\begin{itemize}
\raggedright
  \item \emph{Say it:} \emph{orīgō, orīginis}
  \item \emph{Say it:} \emph{orīgō, orīginis}, once more
  \item \emph{Say it:} a sentence of your own with \emph{orīgō, orīginis}
  \item \emph{Recall:} say the Latin for annals, then the Latin for a founder, then say \emph{orīgō, orīginis} again
\end{itemize}

\begin{tcolorbox}[breakable,colback=teal!4,colframe=teal!35!black,title={Before you move on}]
What does orīgō, orīginis mean? (\textquotedblleft{}An origin, a beginning; ancestry\textquotedblright{}.) Say it once more.
\end{tcolorbox}
\section[vetustās, vetustātis]{vetustās, vetustātis — great age, antiquity}

\label{lesson:LA-C344-vetustas}



\begin{quote}
Before the new one: say the Latin for a founder, then the Latin for an origin.
\end{quote}

\subsection*{You'll want to know: vetustās, vetustātis}
\textbf{vetustās, vetustātis} — \textquotedblleft{}great age, antiquity\textquotedblright{}.

\begin{quote}
Templum vetustāte collāpsum erat.
\end{quote}

\begin{quote}
\emph{The temple had fallen in from age.}
\end{quote}

Feminine, 3rd declension, from \emph{vetus} (taught). \emph{Vetustāte}: through long age.

\begin{cousinweb}
\textbf{vetustās} is built on \emph{vetustus}, \textquotedblleft{}aged\textquotedblright{}, from \emph{vetus}, \textquotedblleft{}old\textquotedblright{} (Lewis \& Short): long existence. English \emph{veteran} belongs to the same family.
\end{cousinweb}

\begin{grammarlens}[title={What you've built}]
One more word for telling, explaining and describing at B1.
\end{grammarlens}

\subsection*{Your turn}
\begin{itemize}
\raggedright
  \item \emph{Say it:} \emph{vetustās, vetustātis}
  \item \emph{Say it:} \emph{vetustās, vetustātis}, once more
  \item \emph{Say it:} a sentence of your own with \emph{vetustās, vetustātis}
  \item \emph{Recall:} say the Latin for a founder, then the Latin for an origin, then say \emph{vetustās, vetustātis} again
\end{itemize}

\begin{tcolorbox}[breakable,colback=teal!4,colframe=teal!35!black,title={Before you move on}]
What does vetustās, vetustātis mean? (\textquotedblleft{}Great age, antiquity\textquotedblright{}.) Say it once more.
\end{tcolorbox}
